\section{案例：5V 外部状态怎样安全进入 3.3V 输入}

“高电平能够被识别”与“引脚能够长期承受该电压”是两个不同条件。设板外设备使用 5V 推挽输出，保证低电平不高于 0.4V、高电平不低于 4.4V；接收芯片工作在 3.3V，输入低阈值为 0.99V、输入高阈值为 2.31V，但绝对最大输入只有 $V_{\mathrm{DD}}+0.3\mathrm{V}$，约 3.6V。低电平同时满足功能与安全要求；高电平虽然远高于 2.31V，能够被判为 1，却超过了 3.6V 的绝对最大值，不能直接连接。

这组数字揭示了数据手册中的两类边界。$V_{\mathrm{IL}}$ 与 $V_{\mathrm{IH}}$ 说明接收器怎样解释电平，属于功能规格；绝对最大额定值说明器件不受损的极限，属于生存边界。绝对最大值不是推荐工作点，也不能因为串联电阻限制了电流，就默认可以持续依赖片内保护结构工作。

\begin{longtable}{L{0.20\linewidth}L{0.24\linewidth}L{0.25\linewidth}L{0.21\linewidth}}
\toprule
输入状态 & 功能判断 & 电气判断 & 结论 \\
\midrule
5V 端输出 0.4V & 低于 0.99V，可靠为 0 & 未超过负向额定范围 & 可作为合法低电平 \\\tablerowrule
5V 端输出 4.4V & 高于 2.31V，可靠为 1 & 高于约 3.6V 的绝对最大值 & 逻辑正确但电气不安全 \\\tablerowrule
接收端未上电 & 不能形成有效内部逻辑 & 高电平可能经钳位路径向 3.3V 电源反灌 & 必须阻断或限制反向供电 \\\tablerowrule
输入处于 0.99--2.31V & 不保证为 0 或 1 & 内部上下拉网络可能同时部分导通 & 不得作为稳定工作状态 \\
\bottomrule
\end{longtable}

\topic{串联电阻能限制异常电流，但不能自动完成电平转换}

若 5V 信号只是一次短暂异常，接收端钳位约在 3.6V，且数据手册允许钳位电流不超过 2mA，那么串联电阻的理论下限为
\[
  R_{\min}=\frac{5.0-3.6}{2\mathrm{mA}}=700\Omega.
\]
取 1k$\Omega$ 时，理想钳位电流约为 1.4mA。这个计算只回答“异常发生时电流是否受限”，没有证明片内钳位适合持续导通，也没有处理接收芯片掉电时的反灌。数据手册若没有明确允许这种工作方式，就应把它视为保护路径，而不是正常数据路径。

对于单向、低速状态信号，可以使用电阻分压。上臂 10k$\Omega$、下臂 20k$\Omega$ 时，5V 被分为约 3.33V；两只电阻的戴维南等效电阻约为 6.67k$\Omega$，还要与输入电容共同决定上升时间，并核对输入漏电在高温下造成的误差。对于高速信号、双向信号、接收端可能掉电的信号，使用明确支持 5V 输入且具备掉电隔离的缓冲器或电平转换器更稳妥。

\begin{pdfdiagram}[x=1cm,y=1cm]
  \node[hw state, minimum width=2.5cm] (source) at (1.4,4.0) {板外 5V 输出\\0.4V 或 4.4V};
  \node[hw control, minimum width=2.7cm] (protect) at (4.6,4.0) {连接器保护\\ESD/浪涌/限流};
  \node[hw compute, minimum width=2.8cm] (translate) at (8.0,4.0) {电平转换\\缓冲或电阻网络};
  \node[hw state, minimum width=2.6cm] (pin) at (11.3,4.0) {3.3V 输入脚\\安全电压范围};
  \node[hw control, minimum width=2.8cm] (threshold) at (11.3,1.2) {输入阈值\\低/未规定/高};
  \node[hw state, minimum width=2.8cm] (bit) at (8.0,1.2) {内部可靠 bit\\明确 0 或 1};
  \node[hw control, minimum width=2.7cm] (fault) at (4.6,1.2) {异常监测\\过压、断线、掉电};
  \node[hw state, minimum width=2.5cm] (safe) at (1.4,1.2) {安全结果\\隔离或默认状态};
  \draw[hw bus] (source) -- (protect);
  \draw[hw bus] (protect) -- (translate);
  \draw[hw bus] (translate) -- (pin);
  \draw[hw bus] (pin) -- (threshold);
  \draw[hw return] (threshold) -- (bit);
  \draw[hw return] (bit) -- (fault);
  \draw[hw return] (fault) -- (safe);
\end{pdfdiagram}
\diagramnote{5V 外部状态进入 3.3V 逻辑的完整边界。上排先限制异常能量并完成电平适配，下排才由输入阈值形成内部 bit；过压、断线或接收端掉电时，故障路径必须把输入隔离或带到规定默认状态，不能让片内钳位长期承担电平转换。}

\topic{掉电顺序决定是否出现反向供电}

接收芯片掉电时，$V_{\mathrm{DD}}$ 变为 0V。若外部 5V 输出仍保持高电平，输入钳位可能把电流送入已经关闭的 3.3V 电源网，使芯片局部处于半上电状态。此时内部复位、I/O 状态和功耗都不再受正常规格保证；电流还可能经其他引脚继续流向板上器件。具备 \term{掉电保护}{power-off protection} 或 \term{部分掉电模式}{partial power-down mode} 的缓冲器，会在自身或接收端未供电时保持高阻，切断这条路径。

调试时应同时测三个点：5V 发送端、转换器输出端和 3.3V 电源轨。正常上电状态下核对高低电平与边沿；接收端断电、发送端保持高电平时，确认 3.3V 电源没有被抬升；快速插拔或 ESD 测试后，确认转换器输出仍回到规定默认状态。仅在正常上电时看到逻辑分析仪显示 0/1，不能证明掉电和异常路径已经闭合。

\begin{longtable}{L{0.20\linewidth}L{0.28\linewidth}L{0.30\linewidth}L{0.12\linewidth}}
\toprule
方案 & 适用条件 & 必须核对的边界 & 评价 \\
\midrule
直接连接 & 接收脚明确标注 5V tolerant & 上电、掉电和所有输入模式均允许 5V & 器件明确支持时最简单 \\\tablerowrule
串联限流 & 短暂异常或手册明确允许钳位电流 & 最大钳位电流、持续时间、反灌路径 & 主要是保护，不默认作为转换 \\\tablerowrule
电阻分压 & 单向、低速、输入漏电小 & 最坏阈值、RC 边沿、功耗与掉电状态 & 简单但速度和方向受限 \\\tablerowrule
专用缓冲/转换器 & 高速、双向或电源时序复杂 & 方向控制、输出使能、掉电高阻与传播延迟 & 接口契约最完整 \\
\bottomrule
\end{longtable}
